\section{Conclusion}
\label{sec:conclusion}

\subsection{Summary}

We have established theoretical and empirical mixing time bounds for
redistricting Markov chains on planar census-tract graphs.

\paragraph{Theoretical.}
For the lazy reversible \recom{} chain on a planar graph with $n$ tracts,
the spectral gap satisfies $\sgap = \Omega(1/n^2)$, implying
$\tmix = O(n^2 \log n)$.
For California ($n = 8{,}057$), this bound is approximately $10^9$ steps.

\paragraph{Empirical.}
In practice, \recom{} achieves $\rhat < 1.1$ (the G.4 convergence
criterion) in 2{,}000--20{,}000 steps, three to four orders of magnitude
faster than the theoretical bound.
The practical mixing time scales approximately as $400k$ (steps per
district), not $O(n^2)$.

\paragraph{ConvergenceSweep.}
The \cs{} algorithm provides a deterministic $O(T \cdot n)$ stopping
criterion that is $\approx 300\times$ faster than certified ensemble
sampling for California.
For plan generation under the DIA statute, \cs{} is unambiguously
preferable.

\paragraph{Division of labor.}
Ensemble methods are irreplaceable for auditing enacted maps and
characterising the feasible plan space.
\cs{} is irreplaceable for generating the statutory plan.
These are complementary roles that should coexist in the DIA framework.

\subsection{Open Problems}

\begin{enumerate}
\item \textbf{Tight mixing time bounds.}
      The $O(n^2 \log n)$ bound is unlikely to be tight; the empirical
      $O(k)$ scaling suggests the true mixing time may be polynomial in
      $k$ rather than $n$.
      A formal proof of this conjecture remains open.

\item \textbf{Non-reversible chains.}
      The spectral gap bound applies to reversible chains.
      The original \recom{} proposal (without Metropolis correction) is
      not reversible, and its mixing time analysis requires different tools.

\item \textbf{Concentration bounds.}
      Formalising the observation that high-compactness plans dominate
      the stationary distribution (explaining the theory-practice gap)
      would require tail bounds on the Polsby-Popper distribution under
      the uniform measure over $\mathcal{P}$.

\item \textbf{Multi-chain convergence.}
      The G.4 diagnostics use $m = 10$ chains.
      The optimal number of chains for efficiency under a fixed step
      budget is an open question.
\end{enumerate}

\subsection{Implications for Redistricting Practice}

The practical upshot of this paper is simple: for states with $k \leq 20$,
run at least 10{,}000 ensemble steps and certify with G.4 diagnostics.
For $k > 20$, run at least $500k$ steps.
Always report $\rhat$, $\ess$, and $\ham(1)$.

For statutory plan generation, use \cs{}.
The MCMC mixing time question is not relevant to plan generation; it
is relevant only to ensemble-based audit methods.
Courts and legislatures should understand this distinction clearly.
